% sec-09-conclusions.tex - Section 9, Conclusions.  DRAFT stage.
% One page maximum.  No figures and no tables: everything in this section has
% already been shown, and the conclusion's job is to state what follows from it.
\section{Conclusions}
\label{sec:conclusions}

Between January 2025 and July 2026 the executive branch built the demand side of
an AI-assisted cancer research program and left the supply side alone. Executive
Order 14355 aimed AI at cancer \cite{eo14355}, the Childhood Cancer Data
Initiative budget doubled \cite{ccdidoubling_nih}, the Genesis Mission and its
NIH arm committed more than \$1.2 billion \cite{eo14363, biogenesis_nih}, the
FDA fielded a generative tool for tasks including clinical protocol review
\cite{crs_ai_healthcare}, and HHS launched a department-wide effort to restore
American leadership in clinical trials alongside Operation TrialBlazer
\cite{trialblazer}. Compute, data, agency review capacity and faster site
activation are all being funded. The document a sponsor must file is not.

\draftinstr{Stage 7: expand this paragraph to name each 2025 to 2026 action once
and connect each to the corresponding capability in Figure 1. One page maximum
for the whole section.}

This paper closes that gap on the sponsor side, and it does so with artifacts
rather than with a proposal. An IND of twelve modules under 21 CFR \S 312 and
\S 812 was assembled in four days \cite{phase1ind}. A Phase 1 protocol at n = 18
behind a Phase 0 gate of at least 1000 simulations across two frameworks, and a
Phase 2 protocol at n = 220 across eight high-volume centers targeting a hazard
ratio of 0.60, were deposited two days apart
\cite{onpremwhippl, phase2proto}. Five bill versions and four companion
documents were deposited in nineteen days \cite{hr9510billv5, clintrust}.
Fourteen funding artifacts followed, carrying a \$1,606,000 small-business ask
against a \$3,500,000 five-year direct program and a virtual trial projected at
\$36,330 against an industry benchmark above \$120,000
\cite{capplan2026, tenapps2026, rfarm27001v2}.

\draftinstr{Stage 7: add one sentence per main section, keyed to sections 3
through 7, each carrying its own strongest quantity. Then state the
incompatibility claim in full.}

The two systems are primarily incompatible, and the incompatibility is not a
matter of preference. A review regime whose best case is seven to eight weeks
and whose typical round is several months cannot absorb an author who deposits
more than thirty artifacts in a year, and a review that arrives after release
cannot correct the object that was released. A production regime whose unit is
the team-month cannot be audited by a commit history it does not keep, and a
figure format that stores pixels cannot be read back as the input to the next
paper. The new system's advantages are persistent rather than incidental:
comprehensive provenance on every artifact, a 1 to 4 day project time scale, AI
peer review during development rather than after it, and machine-readable
figures that are re-read rather than re-drawn.

\draftinstr{Stage 7: close on what the paper is for. The author is seeking near
term funding for the PDAC hybrid clinical trial, and every artifact named above
exists and is deposited under a DOI. Do not overstate: nothing here is an
approved protocol, an active IND, or an enacted statute, and the paper must say
so in its own voice as well as in the cover disclaimer.}
